
\Problem{Donacions}

\Statement

Volem mantenir informació bàsica de les donacions monetàries a un hospital.
En particular, volem guardar quants diners en total ha donat cada persona,
identificada de forma única amb el seu NIF (8 dígits seguits de la lletra majúscula de control).

\medskip

El programa permet fer cinc operacions diferents:
\begin{itemize}
\item
\texttt{N}:
Escriu el nombre de donants precedit de \qq{number:}.
Fixeu-vos que hi pot haver menys donants que donacions,
si hi ha donants repetits.
\item
\texttt{D} $n$ $m$:
S'apunta que la persona amb NIF $n$ ha donat $m$ diners.
Si la persona ja havia fet alguna donació, n'acumula els imports.
\item
\texttt{Q} $n$:
Consulta quants diners ha donat de moment la persona amb NIF $n$.
Cal escriure $-1$ si encara no ha fet cap donació.
\item
\texttt{P}:
Escriu una línia amb tots els NIFs que acaben amb un dígit parell
(ignorant el caràcter de control),
amb els NIFs ordenats alfabèticament.
Cal separar els NIFs amb espais.
\item
\texttt{L}:
Escriu la informació del donant amb el NIF més gran alfabèticament.
Si no n'hi ha cap, escriu \qq{None}.
\end{itemize}


\Input

L'entrada consisteix en diverses operacions segons s'ha indicat anteriorment.
Els NIFs són correctes, i les donacions són nombres naturals.


\Output

Per a cada operació (excepte les donacions),
escriviu la informació demanada.

\Observations

Per aquest exercici, amb un programa sense funcions n'hi ha prou.


\Sample

